% §2 Background — taxonomy of adjacent work in five strands
\subsection{Scope and Taxonomy}

Our previous work analyzed the voice-agent conversation one turn at a time: the latency floor bounds how fast a response can begin \citep{garcia2026latencyfloor}, the cessation floor bounds how fast a response can stop after barge-in \citep{garcia2026cessationfloor}, and the semantic-VAD study supplies a measured Spanish partial-transcript sensing layer \citep{garcia2026asking}.
This paper moves the analysis one level up, from the turn to the complete conversation, treating each call as a trajectory through a grammar of dialogue states whose terminal classes are banking outcomes.
We call the witnessed prefix of a conversation a \emph{partial trajectory} and its eventual terminal class the \emph{outcome}.
The literature that bears on this object sits in five strands that barely touch: agent-trajectory benchmarks, dialogue process structure, partial-observation dialogue management, early outcome prediction, and escalation practice from call-center operations.
We review each strand in turn and close with the delta our grammar adds.

\subsection{Agent-Trajectory Benchmarks: the Full-Episode Paradigm}

The closest benchmark family scores an agent only when its episode is over.
\citet{sierra2024taubench} introduce $\tau$-bench, which evaluates tool-using agents in retail and airline domains against simulated users and domain policies, reporting pass$^k$ over complete trajectories.
The pass$^k$ metric summarizes reliability across $k$ repeated episodes, again measured only at the end.
In $\tau$-bench the trajectory appears as a log to be scored, never as the object of analysis.
WebArena \citep{webarena2023} measures autonomous web agents by comparing final environment state against a goal, inheriting the same full-episode paradigm in a static setting.
$\tau$-Multilingual extends this evaluation style to voice agents in Spanish, Brazilian Portuguese, Hindi, Korean, and Mandarin \citep{tauvoice2026}, confirming that Spanish voice-agent quality measurement is arriving.
Yet its units remain task success and component quality under simulated users, not outcomes predicted from a live, partially observed conversation.
An audit of the paradigm itself sharpens the concern: \citet{falsesuccess2026} characterize false success across 9{,}876 $\tau^2$-bench trajectories, showing that agents assert completion while the environment disagrees.
If terminal labels mislead even in static text environments, conversation-scale evaluation needs outcome evidence that accumulates during the call rather than after it.
The delta of our paper relative to this strand is direct: the strand supplies outcome labels, but no trajectory model, no partial-observation decision rule, and no voice physics.

\subsection{Dialogue Process Structure: Workflows, Quality, and Breakdown}

A second strand studies the structure of conversations as processes.
\citet{salesforce2025workflows} extract dialog workflows from service-agent conversations automatically and evaluate the recovered workflows for consistency and coverage.
The workflow is an artifact distilled after the fact; the conversations that produced it are finished, and no outcome rides on the extraction.
The evaluation tradition is older: \citet{walker1997paradise} established with PARADISE that task success alone underspecifies dialogue quality, which must be traded against dialogue cost.
Our banking KPI protocol inherits this cost axis, but replaces post-hoc quality ratings with per-class costs over trajectory outcomes measured in production telemetry.
Breakdown detection \citep{higashinaka2016breakdown} attributes failure to individual system utterances given their preceding context, a turn-level diagnosis rather than a forecast of where the call is heading.
Closest in spirit, \citet{chang2019derailment} forecast derailment in open-domain social conversation from growing context windows, establishing partial-context prediction as principled.
Their signals, however, are social cues in online forums; there is no task, no tool state, and no action taken on the forecast.
The delta of our paper relative to this strand is that it offers either structure without outcomes or forecasting without decisions, and none of it models voice-agent states such as authentication, tool loops, or barge-ins.

\subsection{Partial-Observation Dialogue Management: the POMDP Lineage}

The classical formal home of "belief over conversation state from noisy partial observations" is the POMDP dialogue-management lineage \citep{young2013pomdp}.
There, a belief over \emph{dialogue state} --- slot values, user goals, dialogue act --- is maintained under ASR noise, and a policy acts on the belief to maximize expected reward.
Our posterior over terminal business outcomes shares this mathematical skeleton but moves the object: terminal banking classes (contained, abandoned, handoff, compliance-escape) are properties of the whole run, not of a turn-level state, and they are only revealed at call teardown.
The POMDP lineage also carries no decision-time guarantee --- its policy optimizes long-run reward, not when a decision may be safely taken --- and it predates LLM tool-calling pipelines, so its state spaces have no tool-event or barge-in channels.
Our grammar keeps the belief-update discipline but attaches it to typed grammar runs, and our stopping theorem adds what the lineage lacks: an explicit bound on when acting on the belief is justified.

\subsection{Early Outcome Prediction: the Nearest Neighbors}

The nearest neighbor is TRACER \citep{nourbakhsh2026tracer}, which formulates task-oriented dialogue failure prediction as partial-dialogue forecasting.
Every visible prefix of a dialogue becomes one supervised instance whose label is the terminal outcome, so a dialogue of length $T$ contributes $T$ instances sharing a single label.
The model is dual-stream: a trajectory stream encodes five hand-built features of belief-state evolution (oscillation score, coverage rate, conflict count, fill velocity, and domain shifts) with a small temporal Transformer, while a text stream encodes serialized belief states with RoBERTa.
The two streams are fused by a multilayer perceptron into a failure probability $\hat{p}_{i,t}$ for dialogue $i$ at turn $t$.
On MultiWOZ~2.4 with oracle belief states, TRACER reaches an AUC of 0.741 and an F1 of 0.463, degrading to 0.617 AUC with generated belief states.
Useful failure signals appear from 25\% of the dialogue visible, and under a false-positive-rate constraint of 0.15 its trigger fires at a mean detection turn of 1.36.
For intervention, a threshold frozen at $\tau = 0.7788$ on the development set converts the score into a recovery trigger, and template-based recovery outperforms a fine-tuned T5 generator on targeting.
Zero-shot transfer is uneven, holding on ABCD (AUC 0.771) but collapsing on SGD (AUC 0.494).
TRACER is candid about what it is: a learned classifier with a tuned threshold, evaluated post hoc on text-chat benchmarks.
It provides no stopping rule, no calibrated confidence, no decision-time guarantee, a single binary outcome, no voice-physics channels such as endpointing errors or barge-ins, and no wiring from outcomes to business KPIs.
Its qualitative failure taxonomy (information incompleteness, contradiction, state instability, and task drift) anticipates, in prose, the drift statistics our grammar makes formal.
EarlyEval \citep{earlyeval2026} shares the prefix-observation insight but applies it offline, predicting episode outcomes early in order to prune expensive agent evaluation rather than to intervene in a running conversation.
At the utterance scale beneath our grammar, X2-Turn \citep{x2turn2026} predicts frame-synchronous turn states jointly with streaming ASR, and ConSol \citep{consol2025} applies sequential probability ratio testing to stop reasoning paths within a single model call.
Neither operates at conversation scale: turn states and reasoning paths are observation signals, not outcome decisions.
The delta of our paper relative to this strand is to keep what it proves (partial transcripts carry outcome signal) and supply what it lacks: a sequential stopping rule with expected decision time bounded at confidence $1-\varepsilon$, multi-class banking outcomes, voice and tool-event observation channels, and intervention latency accounted against published floors.

\subsection{Escalation and Queueing Practice: Decisions Outside the Conversation}

Operations research treats escalation as a queueing decision made outside the conversation.
\citet{queuecontrol2026} study a two-stage human--AI service system and derive when arriving tasks should be routed to automation or to humans, as a function of queue state.
The decision variable is which tasks to automate, computed before any conversation begins and never revised from what is said.
Silent abandonment \citep{silentabandonment2024} estimates customer patience from uncertain contact-center data, modeling abandonment as a property of the queue rather than of the dialogue.
Industrial practice consolidates this outside view: \citet{nubank2026} describe an evaluation-driven framework for support agents at 100-million-user scale, with containment managed as an operational metric rather than a predictable per-call object.
These works supply the cost model we adopt: agent minutes against teller minutes, repeat contacts, and the operational meaning of containment.
What they lack is a conversation-internal signal; nothing in the queueing view can notice, mid-call, that authentication has failed twice and the call is drifting.
Cost-weighted classification rather than accuracy is therefore the right evaluation lens, and our evaluation protocol adopts it.
The delta of our paper relative to this strand is that escalation and abandonment become terminal classes of a trajectory grammar, predicted from the partial transcript and acted on inside the conversation, where any intervention must first clear the cessation floor \citep{garcia2026cessationfloor}.

\subsection{Positioning}

Across the four strands, no work defines a formal state space for voice-agent conversations with terminal banking outcomes, predicts those outcomes from partial observations with a calibrated decision time, or couples the prediction to intervention machinery with measured latency.
Benchmarks score finished episodes, structure mining ignores outcomes, forecasters offer learned scores without guarantees, and operations decisions ignore the conversation.
This intersection, not any single strand, is the contribution space of the paper.
The trajectory grammar, the stopping theorem, and the banking evaluation protocol presented in the remainder of the paper close it.
Voice physics enters as a first-class observation channel, and every intervention inherits the latency discipline of the floors program \citep{garcia2026latencyfloor}.
